%============================================================
%  Bd. 20 - Linkskürzbare Halbgruppen
%============================================================

\documentclass[main.tex]{subfiles}

\ifSubfilesClassLoaded{
  \BandSetupStandalone{B20}
}{
}

\title{Bd. 20 - Linkskürzbare Halbgruppen}
\author{Martin Kunze}
\date{}
\setcounter{file}{20}

\begin{document}

\title{Bd. 20 - Linkskürzbare Halbgruppen}
\author{Martin Kunze}
\date{}
\hypersetup{pageanchor=false}
\maketitle
\hypersetup{pageanchor=true}
\setcounter{tocdepth}{3}
\tableofcontents
\setcounter{file}{20}
\renewcommand{\FormulaBandID}{20}

\chapter{Einführung}

Dieser Band untersucht Halbgruppen mit linkskürzbarer Operation.

\chapter{Linkskürzbare Halbgruppen}

\section{Definition der Linkskürzbarkeit}

\FormulaDefDeltaK[Linkskürzbare Operation]{%
  \mathsf{LK\ddot{u}rz}\!\left(A,\star\right)%
}{LeftCancellativeDef}{%
  \DeltaRow{Mengen}{A\dsep a\dsep b\dsep c}
  \DeltaRow{Linkskürzung}{%
    a,b,c\in A\dsep a\star b=a\star c\vdash b=c}
  \DeltaRow{\textbf{Neues Symbol}}{}
  \DeltaPrem{Linkskürzbarkeit}{\mathsf{LK\ddot{u}rz}\!\left(A,\star\right)}
}

\FormulaAxiomDeltaK[Zugriff auf die Linkskürzung]{%
  \mathsf{LK\ddot{u}rz}\!\left(A,\star\right)\dsep
  a,b,c\in A\dsep a\star b=a\star c
  \vdash b=c%
}{LeftCancellationAxiom}{%
  \DeltaRow{Mengen}{A\dsep a\dsep b\dsep c}
}

\FormulaDefDeltaK[Begriff der linkskürzbaren Halbgruppe]{%
  \LeftCancellativeSemigroupAlg{A}{\star}%
}{LeftCancellativeSemigroupDef}{%
  \DeltaRow{Mengen}{A}
  \DeltaRow{\textbf{Axiome}}{}
  \DeltaPrem{Halbgruppe}{\SemigroupAlg{A}{\star}}
  \DeltaPrem{Linkskürzbarkeit}{%
    \mathsf{LK\ddot{u}rz}\!\left(A,\star\right)}
  \DeltaRow{\textbf{Neues Symbol}}{}
  \DeltaPrem{Linkskürzbare Halbgruppe}{%
    \LeftCancellativeSemigroupAlg{A}{\star}}
}

\FormulaAxiomDeltaK[Zugriff auf die Halbgruppe]{%
  \LeftCancellativeSemigroupAlg{A}{\star}
  \vdash\SemigroupAlg{A}{\star}%
}{LeftCancellativeSemigroupBaseAxiom}{%
  \DeltaRow{Mengen}{A}
}

\FormulaAxiomDeltaK[Zugriff auf die Linkskürzbarkeit]{%
  \LeftCancellativeSemigroupAlg{A}{\star}
  \vdash\mathsf{LK\ddot{u}rz}\!\left(A,\star\right)%
}{LeftCancellativeSemigroupPropertyAxiom}{%
  \DeltaRow{Mengen}{A}
}

\FormulaAxiomDeltaK[Direkte Linkskürzung]{%
  \LeftCancellativeSemigroupAlg{A}{\star}\dsep
  a,b,c\in A\dsep a\star b=a\star c
  \vdash b=c%
}{LeftCancellativeSemigroupCancellationAxiom}{%
  \DeltaRow{Mengen}{A\dsep a\dsep b\dsep c}
}

\section{Beispiele}

\FormulaThmDeltaK[Die Addition ist linkskürzbar]{%
  \mathsf{LK\ddot{u}rz}\!\left(\mathbb{N},+\right)%
}{NaturalAdditionLeftCancellative}{%
  \DeltaRow{Mengen}{\mathbb{N}\dsep a\dsep b\dsep c}
}
\begin{tabproof}
  \proofstep{1}{a\in\mathbb{N}}{\rA}
  \proofstep{2}{b\in\mathbb{N}}{\rA}
  \proofstep{3}{c\in\mathbb{N}}{\rA}
  \proofstep{4}{a+b=a+c}{\rA}
  \proofstep{1,2,3,4}{b=c}{%
    \FormulaRefAuto{PeanoAddLeftCancellation}[thm]{1,2,3,4}}
  \proofstep{}{%
    \mathsf{LK\ddot{u}rz}\!\left(\mathbb{N},+\right)%
  }{\FormulaRefAuto{LeftCancellativeDef}[def]{5}}
\end{tabproof}

\FormulaThmDeltaK[Die Addition bildet eine linkskürzbare Halbgruppe]{%
  \LeftCancellativeSemigroupAlg{\mathbb{N}}{+}%
}{NaturalAdditionLeftCancellativeSemigroup}{%
  \DeltaRow{Mengen}{\mathbb{N}}
}
\begin{tabproof}
  \proofstep{}{\SemigroupAlg{\mathbb{N}}{+}}{%
    \FormulaRefAuto{NaturalAdditionSemigroup}[thm]}
  \proofstep{}{\mathsf{LK\ddot{u}rz}\!\left(\mathbb{N},+\right)}{%
    \FormulaRefAuto{NaturalAdditionLeftCancellative}[thm]}
  \proofstep{}{\LeftCancellativeSemigroupAlg{\mathbb{N}}{+}}{%
    \FormulaRefAuto{LeftCancellativeSemigroupDef}[def]{1,2}}
\end{tabproof}

\paragraph{Einseitiges Beispiel.}
Für $x\star_r y\coloneqq y$ ist jede nichtleere Menge $A$
linkskürzbar. Besitzt $A$ mindestens zwei Elemente, ist diese Halbgruppe
nicht rechtskürzbar.

\section{Morphismen}

\FormulaDefDeltaK[Homomorphismus linkskürzbarer Halbgruppen]{%
  \LeftCancellativeSemigroupHom{f}{A,\star}{B,\diamond}%
}{LeftCancellativeSemigroupHomDef}{%
  \DeltaRow{Mengen}{A\dsep B}
  \DeltaRow{Funktionen}{f\dsep\star\dsep\diamond}
  \DeltaRow{\textbf{Axiome}}{}
  \DeltaPrem{Quellstruktur}{\LeftCancellativeSemigroupAlg{A}{\star}}
  \DeltaPrem{Zielstruktur}{\LeftCancellativeSemigroupAlg{B}{\diamond}}
  \DeltaPrem{Halbgruppenhomomorphismus}{%
    \SemigroupHom{f}{A,\star}{B,\diamond}}
  \DeltaRow{\textbf{Neues Symbol}}{}
  \DeltaPrem{Homomorphismus}{%
    \LeftCancellativeSemigroupHom{f}{A,\star}{B,\diamond}}
}

\FormulaAxiomDeltaK[Quellstruktur]{%
  \LeftCancellativeSemigroupHom{f}{A,\star}{B,\diamond}
  \vdash\LeftCancellativeSemigroupAlg{A}{\star}%
}{LeftCancellativeSemigroupHomSourceAxiom}{%
  \DeltaRow{Mengen}{A\dsep B}
  \DeltaRow{Funktionen}{f\dsep\star\dsep\diamond}
}

\FormulaAxiomDeltaK[Zielstruktur]{%
  \LeftCancellativeSemigroupHom{f}{A,\star}{B,\diamond}
  \vdash\LeftCancellativeSemigroupAlg{B}{\diamond}%
}{LeftCancellativeSemigroupHomTargetAxiom}{%
  \DeltaRow{Mengen}{A\dsep B}
  \DeltaRow{Funktionen}{f\dsep\star\dsep\diamond}
}

\FormulaAxiomDeltaK[Zugriff auf den Halbgruppenhomomorphismus]{%
  \LeftCancellativeSemigroupHom{f}{A,\star}{B,\diamond}
  \vdash\SemigroupHom{f}{A,\star}{B,\diamond}%
}{LeftCancellativeSemigroupHomBaseHomAxiom}{%
  \DeltaRow{Mengen}{A\dsep B}
  \DeltaRow{Funktionen}{f\dsep\star\dsep\diamond}
}

\FormulaDefDeltaK[Isomorphismus linkskürzbarer Halbgruppen]{%
  \LeftCancellativeSemigroupIso{f}{A,\star}{B,\diamond}%
}{LeftCancellativeSemigroupIsoDef}{%
  \DeltaRow{Mengen}{A\dsep B}
  \DeltaRow{Funktionen}{f\dsep\star\dsep\diamond}
  \DeltaRow{\textbf{Axiome}}{}
  \DeltaPrem{Homomorphismus}{%
    \LeftCancellativeSemigroupHom{f}{A,\star}{B,\diamond}}
  \DeltaPrem{Bijektivität}{f\colon A\bij B}
  \DeltaRow{\textbf{Neues Symbol}}{}
  \DeltaPrem{Isomorphismus}{%
    \LeftCancellativeSemigroupIso{f}{A,\star}{B,\diamond}}
}

\FormulaAxiomDeltaK[Homomorphismuszugriff]{%
  \LeftCancellativeSemigroupIso{f}{A,\star}{B,\diamond}
  \vdash\LeftCancellativeSemigroupHom{f}{A,\star}{B,\diamond}%
}{LeftCancellativeSemigroupIsoHomAxiom}{%
  \DeltaRow{Mengen}{A\dsep B}
  \DeltaRow{Funktionen}{f\dsep\star\dsep\diamond}
}

\FormulaAxiomDeltaK[Bijektivitätszugriff]{%
  \LeftCancellativeSemigroupIso{f}{A,\star}{B,\diamond}
  \vdash f\colon A\bij B%
}{LeftCancellativeSemigroupIsoBijectiveAxiom}{%
  \DeltaRow{Mengen}{A\dsep B}
  \DeltaRow{Funktionen}{f\dsep\star\dsep\diamond}
}

\end{document}
